% !TeX program = lualatex
% Definition and sources only; this file is not an LLM solution attempt.
\documentclass[11pt,a4paper]{article}
\usepackage[margin=25mm]{geometry}
\usepackage{fontspec}
\setmainfont{lmroman10-regular.otf}[BoldFont=lmroman10-bold.otf,
 ItalicFont=lmroman10-italic.otf,BoldItalicFont=lmroman10-bolditalic.otf]
\usepackage{amsmath,amssymb,mathtools}
\usepackage{unicode-math}
\setmathfont{latinmodern-math.otf}
\AtBeginDocument{\let\setminus\smallsetminus}
\usepackage{xurl,hyperref}
\hypersetup{colorlinks=true,urlcolor=blue}
\setcounter{secnumdepth}{0}
\setlength{\parindent}{0pt}
\setlength{\parskip}{5pt}
\setlength{\emergencystretch}{3em}
\begin{document}
\section*{Saturated $k$-Sperner Systems of Minimum Size}\label{statement}

\subsection{Definitions and mathematical statement}
For a finite set $X$, write $2^X$ for its power set. A family
$\mathcal F\subseteq2^X$ is $k$-Sperner if it contains no strict chain
$A_0\subsetneq A_1\subsetneq\cdots\subsetneq A_k$ of $k+1$ members.
It is saturated when it is $k$-Sperner and adjoining any missing set
$S\in2^X\setminus\mathcal F$ creates such a chain. Thus saturation is
maximality under inclusion of families, not largest possible cardinality.
Let $\operatorname{sat}_k(n)$ be the smallest cardinality of a saturated
$k$-Sperner family when $|X|=n$.

The question asks whether its exponential growth for sufficiently large
ground sets has an exponent strictly greater than $1/2$. A precise
quantified form of the source's asymptotic assertion is: do there exist
$c>1/2$, a function $n_0:\mathbb N\to\mathbb N$, and nonnegative real
numbers $\eta_k\to0$ such that, for all sufficiently large $k$ and all
$n\ge n_0(k)$,
\[
 \operatorname{sat}_k(n)\ge 2^{(c-\eta_k)k}?
\]
One may increase $n_0(k)$ to ensure $n\ge k$ without altering the problem.
The error term is controlled by $k$ uniformly throughout the permitted
range of $n$, rather than separately for individual ground-set sizes.

The requested improvement is in the exponential rate. A multiplicative
factor polynomial in $k$ in front of $2^{k/2}$ does not by itself give a
constant $c>1/2$. Both the absence of a chain and the creation of a chain
by every possible addition are essential to the minimum-size problem.
\subsection{Short English statement}
For large ground sets, must every saturated family with no chain of length $k+1$ have exponential size with rate strictly above one half?
\subsection{Sources}
\begin{itemize}
\item[S1] ULAM.AI and the UnsolvedMath Contributors, supplied problems.json, catalogue entry, OpenGarden collection. Catalogue identity and source transcription; CC BY 4.0.
\url{https://huggingface.co/datasets/ulamai/UnsolvedMath}.
\item[S2] Open Problem Garden, Saturated $k$-Sperner Systems of Minimum Size. Original problem and discussion; the mathematical formulation below is expanded from the supplied source record.
\url{http://www.openproblemgarden.org/op/saturated_k_sperner_systems_of_minimum_size}.
\end{itemize}
\end{document}
